\chapter{General introduction}

\section{Introduction to linear inverse problems}

The main focus of this section is to present mathematical inverse problems, which are problems where we aim to recover an underlying signal from indirect and noisy measurements. These problems are common in various fields such as image reconstruction, signal processing, and machine learning.
We limit ourselves to linear inverse problems in finite dimensions. This allows us to derive theoretical results and develop efficient algorithms for solving these problems.

Let $(\Omega, \mathcal{F}, \mathbb{P})$ be a probability space and $X: \Omega \to \mathcal{X}$ a random variable representing the true signal. We assume that we have access to measurements $Y: \Omega \to \mathcal{Y}$ that are related to the true signal through a linear operator $A: \mathcal{X} \to \mathcal{Y}$ and corrupted by a noise process $\mathcal{C}$. The measurement model can be written as:
\begin{equation}
    Y | X=x \sim \mathcal{C}(Ax)
    \label{eq:inverse-problem}
\end{equation}

In most setups, we work in high dimensions and the forward operator $A$ is not directly invertible, which makes the problem of recovering $X$ from $Y$ an ill-posed inverse problem.

\begin{definition}
    $A$ is said to be ill-posed if one of the following properties is not satisfied :
\begin{enumerate}
    \item $A$ is injective, i.e. $\forall x_1, x_2 \in \mathcal{X}, Ax_1 = Ax_2 \implies x_1 = x_2$.
    \item $A$ is surjective, i.e. $\forall y \in \mathcal{Y}, \exists x \in \mathcal{X}$ such as $y = Ax$.
    \item $A^{-1}$ is continuous, i.e. $\forall \epsilon > 0, \exists \delta > 0$ such as $\| y_1 - y_2 \| < \delta \implies \| A^{-1}y_1 - A^{-1}y_2 \| < \epsilon$.
\end{enumerate}
\end{definition}

\begin{remark}
    Note that in finite dimention, linear bijections are always continuous and thus the third condition is automatically satisfied.
\end{remark}

If $A$ is ill-posed, then it is obviously not a bijection and its null-space is not reduced to $\{0\}$.
This means that there exist many solutions to the inverse problem for a given measurement, eventhough we're looking for a unique solution.


\begin{definition}{Minimal-norm solution}
    \begin{enumerate}
    \item $x \in \mathcal{X}$ is a least-squares solution if 
    \begin{equation*}
        \| Ax - y \| = \inf_{z \in \mathcal{X}} \| Az - y \|
    \end{equation*}
    \item $x \in \mathcal{X}$ is a minimal-norm solution if it is a least-squares solution and $\| x \| = \inf_{z \in \mathcal{X}, Az = y} \| z \|$.
    \end{enumerate}
\end{definition}

% \begin{definition}{Convergence of solution sequence}
%     Let $(y_\delta)_{\delta > 0}$ be a sequence of measurements such that $\| y_\delta - y \| \leq \delta$.
%     This sequence is said to be convergent with regard to an inverse operator $A^\dagger$ if the sequence of minimal-norm solutions $(x_\delta)_{\delta > 0}$ defined by $x_\delta = A^\dagger y_\delta$ converges to the minimal-norm solution $x^\dagger$ of $y$, i.e. $\lim_{\delta \to 0} \| x_\delta - x^\dagger \| = 0$.

% \end{definition}


\begin{proposition}
    If $A \in \mathcal{L}(\mathcal{X}, \mathcal{Y})$ is a bounded operator and if a least square solution exists, then the minimal-norm solution is unique and is given by $ x ^\dagger = A^\dagger y$ where $A^\dagger$ is the Moore-Penrose pseudo-inverse of $A$. \cite{sherry2025part}
\end{proposition}

\subsection{Regularisation of inverse problems}

If the range of $A$ is not closed, then Moore-Penrose pseudo-inverse may not be well-posed.
Thus, given noisy data $y_\delta$ such that $\| y_\delta - y \| \leq \delta$, the least-squares solution $A^\dagger y_\delta$ may not converge to the minimal-norm solution $x^\dagger$ as $\delta \to 0$.
This is the reason why we need to replace $A^\dagger$ with a family of well-posed and bounded operators $(R_\alpha)_{\alpha > 0}$, called regularisation operators, that approximate $A^\dagger$ as $\alpha \to 0$, hence the following definition.

\begin{definition}{Regularisation}
    A family of continuous operators $(R_\alpha)_{\alpha > 0}$ is called a regularisation of $A^\dagger$ if $\lim_{\alpha \to 0} R_\alpha y = A^\dagger y$ for all $y \in \mathcal{Y}$.
\end{definition}

$\alpha$ is called the regularisation parameter. Setting it to a small value allows us to get a good approximation of the minimal-norm solution, but it also makes the regularisation operator more sensitive to noise. On the other hand, setting $\alpha$ to a large value makes the regularisation operator more robust to noise, but it also leads to a worse approximation of the minimal-norm solution. Therefore, choosing an appropriate value for $\alpha$ is crucial for obtaining good reconstruction results.
This can be seen thanks to the following observation :
\begin{equation*}
    \begin{aligned}
        \| R_\alpha y_\delta - x^\dagger \| &\leq \| R_\alpha y_\delta - R_\alpha y \| + \| R_\alpha y - A^\dagger y \| \\
        & \leq \delta \| R_\alpha \| + \| R_\alpha y - A^\dagger y \|
    \end{aligned}
\end{equation*}
where the first term is some noise amplification term and the second term is some approximation error term. The first term is increasing with $\alpha$ while the second term is decreasing with $\alpha$, which means that there is a trade-off between noise amplification and approximation error when choosing $\alpha$.
This decomposition also indicates that the regularisation parameter $\alpha$ should be chosen according to the noise level $\delta$ in order to achieve a good reconstruction performance.

\subsubsection{Parameter choice rule}

\begin{definition}{Parameter choice rule}
    A parameter choice rule is a function $\alpha: \mathbb{R}^+ \times \mathcal{Y} \to \mathbb{R}^+ \mapsto \alpha(\delta, y_\delta)$.
    There are three parameter choice rule families :
    \begin{itemize}
        \item A priori parameter choice rule : $\alpha$ depends only on $\delta$.
        \item A posteriori parameter choice rule : $\alpha$ depends only on $y_\delta$.
        \item Heuristic parameter choice rule : $\alpha$ depends on both $\delta$ and $y_\delta$.
    \end{itemize}
\end{definition}


In this case, we can still define the minimal-norm solution as the limit of the least-squares solutions for a sequence of measurements that converges to $y$ :

\begin{equation*}
    x^\dagger = \lim_{\delta \to 0} A^\dagger y_\delta
\end{equation*}

\subsection{Spectral regularisation}

\subsubsection{Operator compactness}

\begin{definition}
    An operator $A \in \mathcal{L}(\mathcal{X}, \mathcal{Y})$ is said to be compact if for avery subset $B$ of $\mathcal{X}$ that is bounded, the image closure $\overline{A(B)}$ is compact in $\mathcal{Y}$.
\end{definition}

\begin{theorem}{Singular value decomposition}
    Let $A \in \mathcal{L}(\mathcal{X}, \mathcal{Y})$ be a compact operator. Then there exist :
    
    \begin{itemize}
        \item a sequence of positive real numbers $(\sigma_n)_{n \in \mathbb{N}}$ called singular values of $A$ such that $\sigma_n \to 0$ as $n \to \infty$ and $\sigma_1 \geq \sigma_2 \geq \cdots > 0$.
        \item an orthonormal basis $(x_i)$ of $\text{ker}(A^\perp)$.
        \item an orthonormal basis $(y_i)$ of $\overline{\text{ran}(A)}$
    \end{itemize}
        such that :
    \begin{equation*}
        A x_i = \sigma_i y_i, \quad A^* y_i = \sigma_i x_i
    \end{equation*}
    with $A^*$ being the adjoint of $A$, and we have the following representation of $A$ and $A^*$ :
    \begin{equation*}
        \begin{cases}
            \displaystyle A x = \sum_{n=1}^\infty \sigma_n \langle x, x_n \rangle y_n \\
            \displaystyle A^* y = \sum_{n=1}^\infty \sigma_n \langle y, y_n \rangle x_n
        \end{cases}
    \end{equation*}
\end{theorem}

\begin{theorem}
    Let $A \in \mathcal{L}(\mathcal{X}, \mathcal{Y})$ be a compact operator with singular value decomposition $(\sigma_n, x_n, y_n)_{n \in \mathbb{N}}$ given by the previous theorem. 
    Then the Moore-Penrose pseudo-inverse of $A$ is given by :
    \begin{equation*}
        A^\dagger y = \sum_{n=1}^\infty \frac{1}{\sigma_n} \langle y, y_n \rangle x_n
    \end{equation*}
\end{theorem}

Most linear inverse problems are modelled by compact operators, which means that the singular values of $A$ tend to zero as their index goes to infinity. This is the reason why the Moore-Penrose pseudo-inverse is not well-posed, since it involves dividing by the singular values of $A$ which can be arbitrarily small and very close to zero.

\subsubsection{Truncated singular value decomposition}

A good choice of a regularizer could be the singular value decomposition of $A$ :

\begin{equation*}
    R_\alpha y = \sum_{n=1}^\infty g_\alpha(\sigma_n) \langle y, y_n \rangle x_n
\end{equation*}
with $g_\alpha$ being a family of functions such that $\lim_{\alpha \to 0} g_\alpha(\sigma) = \frac{1}{\sigma}$ for all $\sigma > 0$ and $g_\alpha(0) = 0$ for all $\alpha > 0$.

Then one can choose to set $g_\alpha$ as the truncated singular value decomposition, which is defined as follows :
\begin{equation*}
    g_\alpha(\sigma) = \begin{cases}
        \frac{1}{\sigma} & \text{if } \sigma \geq \alpha \\
        0 & \text{if } \sigma < \alpha
    \end{cases}
\end{equation*}

The regularizer can then be written as :
\begin{equation*}
    R_\alpha y = \sum_{\sigma_n \geq \alpha} \frac{1}{\sigma_n} \langle y, y_n \rangle x_n
\end{equation*}
which is now a finite sum and thus well-posed. However, this regularizer is not very robust to noise since it still involves dividing by the singular values of $A$ which can be close to zero.

\subsubsection{Tikhonov regularisation}

The Tikhonov regularisation also uses the singular value decomposition of $A$ but with a different choice of the function $g_\alpha$ :
\begin{equation*}
    g_\alpha(\sigma) = \frac{\sigma}{\sigma^2 + \alpha}
\end{equation*}

The regularizer can then be written as :
\begin{equation*}
    R_\alpha y = \sum_{n=1}^\infty \frac{\sigma_n}{\sigma_n^2 + \alpha} \langle y, y_n \rangle x_n
\end{equation*}
which is now a well-posed regularizer that is more robust to noise than the truncated singular value decomposition since it does not involve dividing by the singular values of $A$ but rather dividing by a shifted version of the singular values which are always greater than or equal to $\alpha$.



\subsection{Gibbs-Markov Random Fields}

\begin{definition}{Gibbs Fields}
A Gibbs field is a graph where each node is a random variable and the edges represent the relationships between the variables.
A clique is a fully connected subset of nodes in a graph. We denote the set of all cliques as $\mathcal{C}$.
For a given set of random variables $\mathbf{X}$ defined on a Gibbs Field, their joint probability can be written as :

\begin{equation}
    p(\mathbf{X}) = \frac{1}{Z} \prod_{c_i \in \mathcal{C}} \psi_i(c_i)
\end{equation}
where $\psi_i(c_i)$ is an implicit potential function defined on the clique $c_i$ and $Z$ is a normalization constant.
\end{definition}


\begin{definition}{Markov Random Fields (MRFs)}
    A Gibbs field is a graph where each node is a random variable and the edges represent the dependencies between the variables. Two nodes are independant if they are not connected by an edge.
\end{definition}


\begin{theorem}{Hammersley-Clifford Theorem}
    For any choice of positive potential functions $\psi_i$, a Markov Random Field and a Gibbs Field are equivalent with regard to the same graph.
    In other words,
    \begin{itemize}
        \item The Gibbs distribution defined on a Markov Random Field can be expressed as a product of potential functions over the cliques of the graph.
        \item Conversely, any positive distribution that factorizes according to the cliques of a graph can be represented as a Gibbs distribution on that graph.
    \end{itemize} 
\end{theorem}

We will consider that conditions for the Hammersley-Clifford theorem are satisfied, and we will use the terms Gibbs Fields and Markov Random Fields interchangeably.
In the context of image processing, the random variables are the pixel values of the image, and the edges represent the spatial relationships between neighboring pixels. The potential functions $\psi_i$ encode our prior knowledge about the image, such as smoothness or edge preservation.
In practice, the potential functions are often defined on pairs of neighboring pixels, which leads to the following formulation for the regularization term $R(x)$:

\begin{equation}
    R(x) = \sum_{j} \sum_{k \in N(j)} w_{jk} \phi(x_j, x_k)
    \label{eq:mrf}
\end{equation}
where $N(j)$ is the set of neighboring pixels of pixel $j$, $w_{jk}$ is a weight that controls the strength of the interaction between pixels $j$ and $k$, and $\phi$ is a potential function that defines the type of regularization applied to the differences between neighboring pixels.

\paragraph{$l_2$ penalty (Quadratic)}

One of the most basic choice for the potential function $\phi$ is the quadratic function, first proposed by \cite{tikhonov1977solutions}, known as \textbf{Tikhonov regularization}. It encourages global smoothness by penalizing the squared differences between neighboring pixels. The side effect of this regularization is that it can oversmooth the image and blur edges, which are important features in medical imaging as we want to distinguish between different tissues and structures.
\begin{equation}
    \phi(x_j, x_k) = (x_j - x_k)^2
\end{equation}

\paragraph{$l_1$ penalty (Total Variation)}

In order to preserve edges while still encouraging smoothness, the \textbf{Total Variation (TV)} regularization was proposed by \cite{Rudin1992}. It penalizes the absolute differences between neighboring pixels, which allows for sharp transitions in the image while still reducing noise in homogeneous regions.
\begin{equation}
    \phi(x_j, x_k) = |x_j - x_k|
\end{equation}

\paragraph{Huber penalty}

Hybrid penalties combine the advantages of both $l_2$ and $l_1$ regularization by applying a quadratic penalty to small differences and a linear penalty to large differences. This allows for smoothness in homogeneous regions while preserving edges. One example of a hybrid penalty is the Huber function, which is defined as follows:
\begin{equation}
    \phi(x_j, x_k) = \begin{cases}
        \displaystyle \frac{(x_j - x_k)^2}{2 \delta}  & \text{if } |x_j - x_k| \leq \delta \\
        \displaystyle |x_j - x_k| - \frac{\delta}{2} & \text{if } |x_j - x_k| > \delta
    \end{cases}
\end{equation}



\section{Introduction to Positron Emission Tomography (PET)}

\subsection{How PET works}

\subsection{Deterministic reconstruction methods}

\paragraph{Radon transform}

\begin{definition}{Radon Transform}
\label{def:radon_transform}
    The Radon transform is an integral transform that takes a two-dimensional function $f(x,y)$ and computes its line integrals along straight lines parameterized by angle $\phi$ and distance $s$ from the origin.
    Mathematically, the Radon transform is defined as:
    \begin{equation}
        g(s, \phi) = \int \int_{-\infty}^{\infty} f(x,y) \, \delta(s - x \cos \phi - y \sin \phi) \, dx dy
    \end{equation}
\end{definition}

In the context of PET imaging, the Radon transform models the process of acquiring projection data (i.e., sinograms) from the distribution of radiotracer within the body.

\begin{theorem}{Central Slice Theorem}
\label{thm:central_slice}
    The Radon transform is invertible in the frequency domain.
\end{theorem}
\begin{proof}
    We consider the 1D Fourier transform of the projection $g(s,\phi)$ with respect to the variable $s$:
    \begin{equation*}
        G(\omega, \phi) = \int_{-\infty}^{\infty} g(s, \phi) e^{-i \omega s} ds
    \end{equation*}
    Substituting the expression for $g(s,\phi)$ into the Fourier transform, we get:
    \begin{equation*}
        G(\omega, \phi) = \int_{-\infty}^{\infty} \left( \int \int_{-\infty}^{\infty} f(x,y) \, \delta(s - x \cos \phi - y \sin \phi) \, dx dy \right) e^{-i \omega s} ds
    \end{equation*}
    Interchanging the order of integration, we have:
    \begin{equation*}
        G(\omega, \phi) = \int \int_{-\infty}^{\infty} f(x,y) \left( \int_{-\infty}^{\infty} \delta(s - x \cos \phi - y \sin \phi) e^{-i \omega s} ds \right) dx dy
    \end{equation*}
    The inner integral evaluates to $e^{-i \omega (x \cos \phi + y \sin \phi)}$, leading to:
    \begin{equation*}
        G(\omega, \phi) = \int \int_{-\infty}^{\infty} f(x,y) e^{-i \omega (x \cos \phi + y \sin \phi)} dx dy
    \end{equation*}
    This expression represents the 2D Fourier transform of $f(x,y)$ evaluated at the frequency coordinates $(\omega \cos \phi, \omega \sin \phi)$. Thus, we have:
    \begin{equation}
        G(\omega, \phi) = F(\omega \cos \phi, \omega \sin \phi)
    \end{equation}
    where $F(u,v)$ is the 2D Fourier transform of $f(x,y)$. This shows that the 1D Fourier transform of the projection at angle $\phi$ corresponds to a slice of the 2D Fourier transform of the object $f(x,y)$ along the direction defined by $\phi$.
\end{proof}


This theorem comes with limitations in practice as we can only acquire a limited number of projections and the data is noisy. Indeed, going from the frequency domain to the spatial domain requires interpolation, which can introduce artifacts if the sampling is not dense enough.
Hence, it is not possible to perfectly reconstruct the original image using only the Radon transform and its inverse. 

\paragraph{Filtered Back Projection (FBP)}
Another possible reconstruction method is known as Back Projection (BP).
This method consists in smearing back each projection across the image along the direction it was acquired.

\begin{definition}{Back Projection}
    \label{def:back_projection}
    The Back projection operation $b$ is defined as follows:
    \begin{equation}
        b(x, y) = \int_{0}^{\pi} g(s, \phi) d\phi = \int_{0}^{\pi} g(x \cos \phi + y \sin \phi, \phi) d\phi
    \end{equation}
\end{definition}

\begin{proposition}
    The Back Projection operation can be interpreted as a convolution of the original image with a $1/r$ kernel, where $r$ is the distance from the center of the image. This convolution results in a blurring effect, as high-frequency components of the image are attenuated
\end{proposition}

\begin{proof}
    Taking the 2D Fourier transform of the Back Projection operation, we have:
    \begin{align*}
        B(u, v) &= \int_{-\infty}^{\infty} \int_{-\infty}^{\infty} \int_{0}^{\pi} g(s, \phi) e^{-2i \pi (ux + vy)} dx dy d\phi\\
        &= \int_{0}^{\pi} \left( \int_{-\infty}^{\infty} \int_{-\infty}^{\infty} g(s, \phi) e^{-2i \pi (ux + vy)} dx dy \right) d\phi\\
        &= \int_{0}^{\pi} \left( \int_{-\infty}^{\infty} g(s, \phi) e^{-2i \pi s \sqrt{u^2 + v^2}} ds \right) d\phi\\
        &= \int_{0}^{\pi} G(\sqrt{u^2 + v^2}, \phi) d\phi
    \end{align*}
    Substituting the expression for $G(\omega, \phi)$ from the Central Slice Theorem (\ref{thm:central_slice}), we get:
    \begin{equation*}
        B(u,v) = \int_{0}^{\pi} F(\sqrt{u^2 + v^2} \cos \phi, \sqrt{u^2 + v^2} \sin \phi) d\phi
    \end{equation*}
    This integral represents a radial integration of the 2D Fourier transform of the original image $f(x,y)$.
    The result is that the Back Projection operation effectively applies a low-pass filter to the original image, attenuating high-frequency components and resulting in a blurred image.
\end{proof}

To counteract this blurring effect, a filtering step is introduced before the Back Projection, leading to the Filtered Back Projection (FBP) method.

\begin{definition}{Filtered Back Projection}
    The Filtered Back Projection operation consists in applying a high-pass filter to the projections before performing the Back Projection:
    \begin{equation}
        f_{\text{FBP}}(x, y) = \int_{0}^{\pi} \left( g(s, \phi) * h(s) \right) d\phi
    \end{equation}
    where $h(s)$ is a high-pass filter kernel, such as the Ram-Lak filter.
\end{definition}

\subsection{Iterative reconstruction in PET}

\subsubsection{Poisson Bayesian inference}

In PET imaging, the data acquisition process can be modeled as a Poisson process as the number of detected events in each detector pair follows a Poisson distribution.
Thus, the likelihood of observing the measured data $y$ given the image $x$ can be expressed as:
\begin{equation}
    p(y|x) = \prod_{i=1}^{d} \frac{(y_i(x))^{y_i} e^{-y_i(x)}}{y_i!}
\end{equation}
where $y_i(x) = (Ax)_i$ is the expected number of counts in detector pair $i$ given the image $x$. It is modeled as a linear operation of the image through the system matrix $A$. 

The aim is to estimate the image $x$ that maximizes the likelihood of the unobserved data $x$ given the observed data $y$. This is known as the Maximum Likelihood (ML) estimation. We will transform this formulation into a minimization problem by taking the negative log-likelihood:
\begin{equation}
    \hat{x}_{\text{ML}} = \arg \min_{x} \sum_{i=1}^{d} \left( (Ax)_i - y_i \log (Ax)_i\right)
    \label{eq:ml_estimation}
\end{equation}

However, one may not want to find the optimal solution of \eqref{eq:ml_estimation} as $A$ is ill-conditionned. Hence, several strategies such as penalization, iteration truncation or post-filtering can be applied to the output of the algorithm to improve image quality.

\subsubsection{Maximum-Likelihood Expectation-Maximization (MLEM)}

Among iterative algorithms for PET reconstruction, the \textbf{Maximum-Likelihood Expectation-Maximization (MLEM)} proposed by \cite{Shepp1982} is one of the most widely used. It is an iterative method to find the ML estimate \eqref{eq:ml_estimation} of the image $x$ given the measured data $y$.

The MLEM algorithm is guaranteed to converge to a global minimum, providing an estimate of the image $x$ that is consistent with the measured data $y$ under the Poisson noise model \cite{8467338}.
However, it can be slow to converge and can produce noisy images if the number of iterations is too high.
In order to address this, iterations can be stopped early, or output can be post-processed or post-filtered.

\begin{algorithm}
\caption{Maximum-Likelihood Expectation-Maximization (MLEM)}
\begin{algorithmic}[1]
\State \textbf{Input:} Measured data $y$, system matrix $A$, initial estimate $x^{(0)}$
\State \textbf{Output:} Estimated image $x^{(k)}$
\For{$k = 0$ to $K-1$}
    \State \textbf{E-step:} Compute the expected number of counts
    \begin{equation*}
        \hat{y}_i^{(k)} = (Ax^{(k)})_i
    \end{equation*}
    \State \textbf{M-step:} Update the image estimate
    \begin{equation*}
        x_j^{(k+1)} = x_j^{(k)} \frac{\sum_{i=1}^{d} A_{ij} \frac{y_i}{\hat{y}_i^{(k)}}}{\sum_{i=1}^{d} A_{ij}}
    \end{equation*}
\EndFor
\State \textbf{Return} $x^{(K)}$
\end{algorithmic}
\label{alg:mlem}
\end{algorithm}

% \begin{proposition}
%     The MLEM algorithm converges to a local maximum of the likelihood function, providing an estimate of the image $x$ that is consistent with the measured data $y$ under the Poisson noise model.
% \end{proposition}

% \begin{proof}
%     First we compute the complete data ; $z_{ij}$ is the number of photons emitted from voxel $j$ and detected by detector pair $i$. Hence, $y_i$ = $\sum_{j} z_{ij}$ and $\bar{z_{ij}} = a_{ij} x_j$.
%     The complete data likelihood is given by:
%     \begin{align*}
%         &zp(z|x) = \prod_{i=1}^{d} \prod_{j=1}^{n} \frac{(a_{ij} x_j)^{z_{ij}}}{z_{ij}!} e^{-a_{ij} x_j} \\
%         & \Rightarrow \log p(z|x) = \sum_{i=1}^{d} \sum_{j=1}^{n} \left( z_{ij} \log (a_{ij} x_j) - a_{ij} x_j + K \right) 
%     \end{align*}
%     Then we compute its derivative with respect to $x_j$ and set it to zero to find the maximum:
%     \begin{equation*}
%         \frac{\partial \log p(z|x)}{\partial x_j} = \sum_{i=1}^{d} \left( \frac{z_{ij}}{x_j} - a_{ij} \right) = 0
%     \end{equation*}
%     This leads to the update rule:
%     \begin{equation*}
%         x_j = \frac{\sum_{i=1}^{d} z_{ij}}{\sum_{i=1}^{d} a_{ij}}
%     \end{equation*}
%     Since $z_{ij}$ is unobserved, we replace it with its expectation given the current estimate $x^{(k)}$:
%     \begin{equation*}
%         \hat{z}_{ij}^{(k)} = \mathbb{E}[z_{ij} | y, x^{(k)}] = a_{ij} x_j^{(k)} \frac{y_i}{(Ax^{(k)})_i}
%     \end{equation*}
%     Substituting this into the update rule gives the M-step of the MLEM algorithm:
%     \begin{equation*}
%         x_j^{(k+1)} = x_j^{(k)} \frac{\sum_{i=1}^{d} a_{ij} \frac{y_i}{(Ax^{(k)})_i}}{\sum_{i=1}^{d} a_{ij}}
%     \end{equation*}
%     This iterative process continues until convergence, leading to an estimate of $x$ that maximizes the likelihood function.
% \end{proof}

\subsubsection{Ordered Subset Expectation Maximization (OSEM)}

In order to speed up the convergence of the MLEM algoritm,
the \textbf{Ordered Subset Expectation Maximization (OSEM)} algorithm was proposed by \cite{Hudson1994}. It is a variant of the MLEM algorithm that divides the measured sinogram $y$ into subsets $\left\{y_1, \hdots, y_Q\right\}$ and performs the E-step and M-step from \ref{alg:mlem} on each subset sequentially (see algorithm \ref{alg:osem}).

Even if the algorithm is not designed to minimized a well-defined cost funcion, the latter converges to an MLEM image in the case of noiseless data, but this might not be true in practice and the solution might not even converge but cycle between similar but distinct images \cite{6152654}.
However, it is faster to get quality results compared to the MLEM algorithm and can be used in practice with a large number of subsets without significant loss of image quality.
Post-processing or post-filtering can also be applied to the output of the OSEM algorithm to further improve image quality.

\begin{algorithm}
\caption{Ordered Subset Expectation Maximization (OSEM)}
\begin{algorithmic}[1]
\State \textbf{Input:} Measured data $y$, system matrix $A$, initial estimate $x^{(0)}$, number of subsets $Q$
\State \textbf{Output:} Estimated image $x^{(k)}$
\For{$k = 0$ to $K-1$}
    \For{$q = 1$ to $Q$}
        \State \textbf{E-step:} Compute the expected number of counts for subset $q$
        \begin{equation*}
            \hat{y}_{i}^{(k,q)} = (Ax^{(k)})_i \quad \text{for } i \in \text{subset } q
        \end{equation*}
        \State \textbf{M-step:} Update the image estimate using subset $q$
        \begin{equation*}
            x_j^{(k,q)} = x_j^{(k,q-1)} \frac{\sum_{i \in \text{subset } q} A_{ij} \frac
            {y_i}{\hat{y}_i^{(k,q)}}}{\sum_{i \in \text{subset } q} A_{ij}}
        \end{equation*}
    \EndFor
\EndFor
\State \textbf{Return} $x^{(K)}$
\end{algorithmic}
\label{alg:osem}
\end{algorithm}

\subsubsection{Penalized reconstruction in PET}

This reconstruction problem can also be formulatted as a regularized optimization by introducing some prior knowledge on the image $x$.
Indeed, Bayes's rule gives us the following expression for the posterior distribution of $x$ given $y$:
\begin{equation}
    p(x|y) = \frac{p(y|x) p(x)}{p(y)}
\end{equation}
where $p(x)$ is the prior distribution of the image $x$ and $p(y)$ is the marginal likelihood of the data $y$.
Once again, by taking the negative log-likelihood of $p(x|y)$, we can transform the problem into a penalized likelihood minimization problem:
\begin{equation}
    \hat{x}_{\text{MAP}} = \arg \min_{x} \left( -\log p(y|x) - \log p(x) \right)
\end{equation}
where $-\log p(y|x)$ is the data fidelity term and $-\log p(x)$ is the regularization term that encodes our prior knowledge on the image $x$.
This prior can be expressed in the following exponential form :

\begin{equation}
    p(x) \propto e^{-\beta R(x)}
\end{equation}
where $R(x)$ is a regularization function and $\beta$ is a regularization parameter that controls the trade-off between data fidelity and regularization. The optimisation hence becomes:

\begin{align}
    \hat{x}_{\text{MAP}} &= \arg \min_{x} \left( -\log p(y|x) + \beta R(x) \right) \\
    &= \arg \min_{x} \left( \sum_{i=1}^{d} \left( (Ax)_i - y_i \log (Ax)_i\right) + \beta R(x) \right)
    \label{eq:map_estimation}
\end{align}

The penalized likelihood approach allows us to incorporate various types of prior knowledge about the image, such as smoothness, sparsity, or anatomical information.
The choice of the regularization function $R(x)$ must alleviate the ill-posedness of the reconstruction problem while preserving important features of the image, such as edges and textures.
The following paragraph presents some commonly used regularization functions in PET image reconstruction.


\subparagraph{Relative Difference penalty}

The relative differences penalty is a concave prior that has been shown to be effective in PET imaging \cite{Nuyts2002}.

\begin{equation}
    \phi(x_j, x_k) = \frac{(x_j - x_k)^2}{x_j + x_k + \gamma |x_j - x_k|}
\end{equation}
where the $\gamma$ parameter controls which relative diferences are considered as large.


This prior information can be used in a similar iterative algorithm as MLEM or OSEM. It is called the \textbf{Maximum A Posteriori Expectation-Maximization (MAPEM)} algorithm, and it consists in replacing the M-step of the MLEM algorithm with a step that takes into account the chosen prior $R(x)$ \cite{Green1990}.

\begin{algorithm}
\caption{Maximum A Posteriori Expectation-Maximization (MAPEM)}
\begin{algorithmic}[1]
\State \textbf{Input:} Measured data $y$, system matrix $A$, initial estimate $x^{(0)}$, regularization parameter $\beta$
\State \textbf{Output:} Estimated image $x^{(k)}$
\For{$k = 0$ to $K-1$}
    \State \textbf{E-step:} Compute the expected number of counts
    \begin{equation*}
        \hat{y}_i^{(k)} = (Ax^{(k)})_i
    \end{equation*}
    \State \textbf{M-step:} Update the image estimate by minimizing the penalized likelihood function
    \begin{equation*}
        x_j^{(k+1)} = \frac{x_j^{(k)}}{\displaystyle \sum_{i=1}^{d} A_{ij} + \beta \left.\frac{\partial R(x)}{\partial x_j}\right|_{x=x^{(k)}}}  \sum_{i=1}^{d} A_{ij} \frac{y_i}{\hat{y}_i^{(k)}}
    \end{equation*}
\EndFor
\State \textbf{Return} $x^{(K)}$
\end{algorithmic}
\label{alg:mapem}
\end{algorithm}